% GENERATED FILE. Edit canonical lessons, then run npm run generate:books.
% canonical-source-hash: fnv1a64:214f64c05dc10c28
% canonical-lessons: TA-C436-timinkalam, TA-C436-cura, TA-C436-nantu, TA-C436-cippi, TA-C436-pavalam

\chapter{Life in the Sea}
\label{ch:ta-life-in-the-sea}
\hlchaptermodality{436}

\begin{chapteropening}
\textbf{By the end of this chapter:} \emph{I can talk about whales, sharks, crabs, seashells and coral.}

Complete the last lesson of chapter 436: I can talk about whales, sharks, crabs, seashells and coral.
\end{chapteropening}

\section[\texorpdfstring{\ta{திமிங்கலம்} (timiṅkalam)}{திமிங்கலம் (timiṅkalam)}]{\ta{திமிங்கலம்} (timiṅkalam) — a whale}

\label{lesson:TA-C436-timinkalam}



\begin{quote}
Before the new one: say the Tamil for a salt pan, then the Tamil for a pebble.
\end{quote}

\subsection*{You'll want to know: \ta{திமிங்கலம்}}
\textbf{\ta{திமிங்கலம்}} — \emph{timiṅkalam} — \textquotedblleft{}a whale\textquotedblright{}.

\begin{quote}
\ta{படகிலிருந்து} \ta{ஒரு} \ta{திமிங்கலத்தைப்} \ta{பார்த்தோம்}.
\end{quote}

\begin{quote}
\emph{paṭakiliruntu oru timiṅkalattaip pārttōm.}
\end{quote}

\begin{quote}
We saw a whale from the boat.
\end{quote}

A -\ta{அம்} noun: \ta{திமிங்கலத்தை}, the whale (object).

\begin{cousinweb}
\textbf{\ta{திமிங்கலம்}} (\emph{timiṅkalam}) is borrowed from Sanskrit \emph{timiṃgila}, \textquotedblleft{}a large fabulous fish\textquotedblright{} (Monier-Williams).
\end{cousinweb}

\begin{grammarlens}[title={What you've built}]
One more word for getting around.
\end{grammarlens}

\subsection*{Your turn}
\begin{itemize}
\raggedright
  \item \emph{Say it:} \emph{timiṅkalam}
  \item \emph{Say it:} \emph{timiṅkalam}, once more
  \item \emph{Say it:} \emph{paṭakiliruntu oru timiṅkalattaip pārttōm.} — We saw a whale from the boat.
  \item \emph{Recall:} say the Tamil for a salt pan, then the Tamil for a pebble, then say \emph{timiṅkalam} again
\end{itemize}

\begin{tcolorbox}[breakable,colback=teal!4,colframe=teal!35!black,title={Before you move on}]
What does \ta{திமிங்கலம்} mean? (\textquotedblleft{}A whale\textquotedblright{}.) Say it once more.
\end{tcolorbox}
\section[\texorpdfstring{\ta{சுறா} (cuṟā)}{சுறா (cuṟā)}]{\ta{சுறா} (cuṟā) — a shark}

\label{lesson:TA-C436-cura}



\begin{quote}
Before the new one: say the Tamil for a pebble, then the Tamil for a whale.
\end{quote}

\subsection*{You'll want to know: \ta{சுறா}}
\textbf{\ta{சுறா}} — \emph{cuṟā} — \textquotedblleft{}a shark\textquotedblright{}.

\begin{quote}
\ta{மீனவர்களின்} \ta{வலையில்} \ta{ஒரு} \ta{சிறிய} \ta{சுறா} \ta{சிக்கியது}.
\end{quote}

\begin{quote}
\emph{mīṉavarkaḷiṉ valaiyil oru ciṟiya cuṟā cikkiyatu.}
\end{quote}

\begin{quote}
A small shark was caught in the fishermen's net.
\end{quote}

\ta{சுறாமீன்} is the full name.

\begin{cousinweb}
\textbf{\ta{சுறா}} (\emph{cuṟā}) is native Dravidian (DEDR 2710), \textquotedblleft{}a shark, a swordfish\textquotedblright{}.
\end{cousinweb}

\begin{grammarlens}[title={What you've built}]
One more word for getting around.
\end{grammarlens}

\subsection*{Your turn}
\begin{itemize}
\raggedright
  \item \emph{Say it:} \emph{cuṟā}
  \item \emph{Say it:} \emph{cuṟā}, once more
  \item \emph{Say it:} \emph{mīṉavarkaḷiṉ valaiyil oru ciṟiya cuṟā cikkiyatu.} — A small shark was caught in the fishermen's net.
  \item \emph{Recall:} say the Tamil for a pebble, then the Tamil for a whale, then say \emph{cuṟā} again
\end{itemize}

\begin{tcolorbox}[breakable,colback=teal!4,colframe=teal!35!black,title={Before you move on}]
What does \ta{சுறா} mean? (\textquotedblleft{}A shark\textquotedblright{}.) Say it once more.
\end{tcolorbox}
\section[\texorpdfstring{\ta{நண்டு} (naṇṭu)}{நண்டு (naṇṭu)}]{\ta{நண்டு} (naṇṭu) — a crab}

\label{lesson:TA-C436-nantu}



\begin{quote}
Before the new one: say the Tamil for a whale, then the Tamil for a shark.
\end{quote}

\subsection*{You'll want to know: \ta{நண்டு}}
\textbf{\ta{நண்டு}} — \emph{naṇṭu} — \textquotedblleft{}a crab\textquotedblright{}.

\begin{quote}
\ta{மணலில்} \ta{நண்டுகள்} \ta{ஓடின}.
\end{quote}

\begin{quote}
\emph{maṇalil naṇṭukaḷ ōṭiṉa.}
\end{quote}

\begin{quote}
Crabs ran across the sand.
\end{quote}

The plural is \ta{நண்டுகள்}.

\begin{cousinweb}
\textbf{\ta{நண்டு}} (\emph{naṇṭu}) is native Dravidian (DEDR 2901), \textquotedblleft{}a crab, a lobster\textquotedblright{}.
\end{cousinweb}

\begin{grammarlens}[title={What you've built}]
One more word for getting around.
\end{grammarlens}

\subsection*{Your turn}
\begin{itemize}
\raggedright
  \item \emph{Say it:} \emph{naṇṭu}
  \item \emph{Say it:} \emph{naṇṭu}, once more
  \item \emph{Say it:} \emph{maṇalil naṇṭukaḷ ōṭiṉa.} — Crabs ran across the sand.
  \item \emph{Recall:} say the Tamil for a whale, then the Tamil for a shark, then say \emph{naṇṭu} again
\end{itemize}

\begin{tcolorbox}[breakable,colback=teal!4,colframe=teal!35!black,title={Before you move on}]
What does \ta{நண்டு} mean? (\textquotedblleft{}A crab\textquotedblright{}.) Say it once more.
\end{tcolorbox}
\section[\texorpdfstring{\ta{சிப்பி} (cippi)}{சிப்பி (cippi)}]{\ta{சிப்பி} (cippi) — a seashell; an oyster}

\label{lesson:TA-C436-cippi}



\begin{quote}
Before the new one: say the Tamil for a shark, then the Tamil for a crab.
\end{quote}

\subsection*{You'll want to know: \ta{சிப்பி}}
\textbf{\ta{சிப்பி}} — \emph{cippi} — \textquotedblleft{}a seashell; an oyster\textquotedblright{}.

\begin{quote}
\ta{குழந்தைகள்} \ta{கடற்கரையில்} \ta{சிப்பிகளைச்} \ta{சேகரித்தார்கள்}.
\end{quote}

\begin{quote}
\emph{kuḻantaikaḷ kaṭaṟkaraiyil cippikaḷaic cēkarittārkaḷ.}
\end{quote}

\begin{quote}
The children collected seashells on the beach.
\end{quote}

\ta{முத்துச்} \ta{சிப்பி} is a pearl oyster.

\begin{cousinweb}
\textbf{\ta{சிப்பி}} (\emph{cippi}) is native Dravidian (DEDR 2535), \textquotedblleft{}a shell, a shellfish\textquotedblright{}.
\end{cousinweb}

\begin{grammarlens}[title={What you've built}]
One more word for getting around.
\end{grammarlens}

\subsection*{Your turn}
\begin{itemize}
\raggedright
  \item \emph{Say it:} \emph{cippi}
  \item \emph{Say it:} \emph{cippi}, once more
  \item \emph{Say it:} \emph{kuḻantaikaḷ kaṭaṟkaraiyil cippikaḷaic cēkarittārkaḷ.} — The children collected seashells on the beach.
  \item \emph{Recall:} say the Tamil for a shark, then the Tamil for a crab, then say \emph{cippi} again
\end{itemize}

\begin{tcolorbox}[breakable,colback=teal!4,colframe=teal!35!black,title={Before you move on}]
What does \ta{சிப்பி} mean? (\textquotedblleft{}A seashell; an oyster\textquotedblright{}.) Say it once more.
\end{tcolorbox}
\section[\texorpdfstring{\ta{பவளம்} (pavaḷam)}{பவளம் (pavaḷam)}]{\ta{பவளம்} (pavaḷam) — coral}

\label{lesson:TA-C436-pavalam}



\begin{quote}
Before the new one: say the Tamil for a crab, then the Tamil for a seashell.
\end{quote}

\subsection*{You'll want to know: \ta{பவளம்}}
\textbf{\ta{பவளம்}} — \emph{pavaḷam} — \textquotedblleft{}coral\textquotedblright{}.

\begin{quote}
\ta{கடலுக்கு} \ta{அடியில்} \ta{பவளப்} \ta{பாறைகள்} \ta{இருக்கின்றன}.
\end{quote}

\begin{quote}
\emph{kaṭalukku aṭiyil pavaḷap pāṟaikaḷ irukkiṉṟaṉa.}
\end{quote}

\begin{quote}
There are coral reefs under the sea.
\end{quote}

\ta{பவளப்} \ta{பாறை} is a coral reef.

\begin{cousinweb}
\textbf{\ta{பவளம்}} (\emph{pavaḷam}) is native Dravidian (DEDR 3998), \textquotedblleft{}red coral\textquotedblright{}.
\end{cousinweb}

\begin{grammarlens}[title={What you've built}]
One more word for getting around.
\end{grammarlens}

\subsection*{Your turn}
\begin{itemize}
\raggedright
  \item \emph{Say it:} \emph{pavaḷam}
  \item \emph{Say it:} \emph{pavaḷam}, once more
  \item \emph{Say it:} \emph{kaṭalukku aṭiyil pavaḷap pāṟaikaḷ irukkiṉṟaṉa.} — There are coral reefs under the sea.
  \item \emph{Recall:} say the Tamil for a crab, then the Tamil for a seashell, then say \emph{pavaḷam} again
\end{itemize}

\begin{tcolorbox}[breakable,colback=teal!4,colframe=teal!35!black,title={Before you move on}]
What does \ta{பவளம்} mean? (\textquotedblleft{}Coral\textquotedblright{}.) Say it once more.
\end{tcolorbox}
